\chapter{Tokenization}

Language models don't read raw text---they process sequences of integers.
Tokenization is the step that converts a string like ``The cat sat on the
mat'' into a list of numbers like $[464, 3797, 5832, 319, 262, 2603]$.
Each integer is an index into an \emph{embedding table}, and the model
learns to operate on these indices rather than characters or bytes.
Choosing how to break text into tokens is one of the first and most
consequential design decisions when building a language model.

This chapter covers the full tokenization pipeline: the spectrum of
approaches (character-level, word-level, subword), a simple
whitespace tokenizer, pre-tokenization, byte-level BPE, the BPE algorithm,
and how to construct next-token prediction datasets from tokenized
corpora.

\section{The Tokenization Spectrum}

Consider the sentence ``She walked quickly.'' The three main
tokenization strategies produce very different outputs:

\begin{center}
\begin{tabular}{lll}
\toprule
\textbf{Strategy} & \textbf{Tokens} & \textbf{IDs (example)} \\
\midrule
Character-level & S, h, e, \textvisiblespace, w, a, l, k, e, d, \ldots & [33, 8, 5, 0, 29, 3, 14, 15, 5, 7, \ldots] \\
Word-level & She, walked, quickly, . & [184, 4052, 2187, 12] \\
Subword (BPE) & She, walk, ed, quickly, . & [184, 3082, 753, 2187, 12] \\
\bottomrule
\end{tabular}
\end{center}

\begin{keyconcept}
  There is a fundamental tradeoff in tokenizer design:
  \begin{itemize}
    \item \textbf{Character-level}: tiny vocabulary ($\sim$256 entries, one per
          ASCII character). Produces very long sequences, making the model's
          attention cost explode. No semantic grouping---the letter ``e'' in
          ``the'' and ``believe'' carries no word-level meaning.
    \item \textbf{Word-level}: short sequences, but the vocabulary is
          unbounded (new words, typos, and morphological variants all need
          entries). Out-of-vocabulary (OOV) words must be mapped to a single
          \texttt{<|unk|>} token, losing information. A vocabulary of 500{,}000+
          words is common for large corpora.
    \item \textbf{Subword (BPE)}: strikes a balance. Common words stay whole
          (``the'' $\rightarrow$ one token), while rare or unseen words decompose
          into known subword units (``unbelievable'' $\rightarrow$
          ``un'', ``believ'', ``able''). Vocabulary size is bounded
          (typically 30k--100k), and every word can be represented. Most modern
          LLMs use this approach.
  \end{itemize}
\end{keyconcept}

\section{SimpleTokenizer}

The simplest tokenizer splits text on whitespace and looks up each token
in a vocabulary. This whitespace splitting is the simplest form of
\emph{pre-tokenization}---the step that breaks raw text into chunks before
subword algorithms operate on them. More sophisticated tokenizers use a
regex pattern instead (see Section~\ref{sec:pretok}).

\begin{pseudocode}[SimpleTokenizer: encode and decode]
class SimpleTokenizer:
    def __init__(self, vocab: dict[str, int]):
        self.vocab = vocab                          # token string -> ID
        self.reverse_vocab = {v: k for k, v in vocab.items()}  # ID -> token string
        self.unk_id = vocab.get("<|unk|>", 0)
        self.eos_id = vocab.get("<|eos|>", 1)

    def encode(self, text: str) -> list[int]:
        tokens = text.split()                      # split on whitespace
        ids = []
        for token in tokens:
            ids.append(self.vocab.get(token, self.unk_id))
        return ids                                 # e.g. [464, 3797, 5832]

    def decode(self, ids: list[int]) -> str:
        tokens = [self.reverse_vocab.get(i, "<|unk|>") for i in ids]
        return " ".join(tokens)                    # concatenate with spaces
\end{pseudocode}

\begin{keyconcept}
  \textbf{Decoding} reverses the encoding process: map each ID back to its
  token string and concatenate. For the SimpleTokenizer, this is
  straightforward---just join with spaces. For BPE, decoding is slightly
  more involved: subword tokens are concatenated directly (without spaces),
  but GPT-2 uses a special character \texttt{\u0120} (derived from the
  byte-to-unicode mapping) to mark tokens that begin with a space. During
  decoding, \texttt{\u0120} is converted back to a space, while tokens
  without it are joined directly.

  Special tokens serve important roles:
  \texttt{<|unk|>} replaces any word not in the vocabulary, giving the model
  a way to represent unseen words rather than crashing;
  \texttt{<|eos|>} marks document boundaries so the model learns where
  sequences end and a new context begins.
\end{keyconcept}

\section{Pre-Tokenization\label{sec:pretok}}

Before BPE can merge pairs, we must first split the raw text into
\emph{words} (or word-like chunks). This pre-tokenization step prevents
BPE from merging characters across word boundaries---we don't want ``he''
in ``the cat'' to merge into the same token as the word ``he.''  The
SimpleTokenizer above used whitespace splitting; real tokenizers use a
more refined regex pattern. GPT-2's pre-tokenizer, for example, splits
text so that contractions, numbers, and punctuation become separate
chunks:

\begin{pseudocode}[Pre-Tokenization Regex (GPT-2)]
import regex                           # not 're' -- see note below

GPT2_PATTERN = r"""'s|'t|'re|'ve|'m|'ll|'d| ?\p{L}+| ?\p{N}+| ?[^\s\p{L}\p{N}]+|\s+(?!\S)|\s+"""

def pretokenize(text: str) -> list[str]:
    return regex.findall(GPT2_PATTERN, text)
# Example: "I'm happy." -> ["I", "'m", " happy", "."]
\end{pseudocode}

\noindent Note that the leading space is \emph{part of the token}
\texttt{"~happy"}: GPT-2's regex attaches spaces to the following word
(rather than making them separate tokens). This preserves word boundaries
during decoding---when we concatenate tokens back, ``~'' marks the start
of a word that was preceded by a space, so we can reconstruct the original
spacing exactly.

Python's built-in \texttt{re} module does not support Unicode property
escapes like \texttt{\textbackslash p\{L\}}; you need the third-party
\texttt{regex} package (\texttt{pip install regex}).

Pre-tokenization determines the \emph{units} that BPE operates on. BPE
merges happen \emph{within} each pre-tokenized chunk, never across chunk
boundaries. This keeps tokens semantically coherent and speeds up training
significantly (the algorithm only counts pairs within each chunk, not
across the entire raw character stream).

\section{Byte-Level BPE}

The BPE variant used by GPT-2 and most modern tokenizers is
\emph{byte-level BPE}. Instead of starting from Unicode characters (which
number in the hundreds of thousands), it starts from the 256 possible byte
values. Any input string---whether English, Chinese, emoji, or binary
data---can be represented as a sequence of bytes, so the initial
vocabulary is exactly 256 entries and every possible input is
representable without an \texttt{<|unk|>} token.

In practice, GPT-2 maps each byte value to a printable character (e.g.,
byte 32 becomes the space character `` ''). This mapping is called the
\emph{byte-to-unicode} table. BPE merges then operate on these byte-level
characters exactly as described in the algorithm below. The resulting
vocabulary consists of the 256 byte-level base tokens plus the learned
merge tokens plus special tokens, yielding GPT-2's total vocabulary size
of 50,257.

\section{Byte-Pair Encoding (BPE)}

Training proceeds in rounds. Each round, we count every adjacent pair
across the entire corpus, pick the most frequent pair, merge it into a
single token, and record the merge rule. The vocabulary grows by one
entry per merge.

\begin{keyconcept}[Worked Example: BPE Training]
  Suppose our (tiny) corpus contains the word ``lowest'' with frequency 5
  and the word ``low'' with frequency 2. We start by splitting each word
  into characters:

  \begin{center}
  \begin{tabular}{lll}
  \toprule
  \textbf{Word} & \textbf{Freq} & \textbf{Character Split} \\
  \midrule
  low & 2 & l\;\, o\;\, w \\
  lowest & 5 & l\;\, o\;\, w\;\, e\;\, s\;\, t \\
  \bottomrule
  \end{tabular}
  \end{center}

  \textbf{Round 1}: Count adjacent pairs weighted by frequency.
  The pair (l, o) appears $2 + 5 = 7$ times; (o, w) appears $2 + 5 = 7$ times;
  (w, e) appears 5 times; (e, s) appears 5 times; (s, t) appears 5 times.
  Ties are broken arbitrarily---in practice, the tie-breaking
  strategy can affect tokenizer quality, so implementations use
  deterministic rules (e.g., alphabetical order of the pair) to ensure
  reproducibility. Here, suppose we pick (l, o). We merge into ``lo'':

  \begin{center}
  \begin{tabular}{lll}
  \toprule
  \textbf{Word} & \textbf{Freq} & \textbf{After Merge} \\
  \midrule
  low & 2 & lo\;\, w \\
  lowest & 5 & lo\;\, w\;\, e\;\, s\;\, t \\
  \bottomrule
  \end{tabular}
  \end{center}

  \textbf{Round 2}: Now (lo, w) appears 7 times, (w, e) appears 5 times,
  (e, s) appears 5 times, (s, t) appears 5 times. Merge (lo, w) into
  ``low'':

  \begin{center}
  \begin{tabular}{lll}
  \toprule
  \textbf{Word} & \textbf{Freq} & \textbf{After Merge} \\
  \midrule
  low & 2 & low \\
  lowest & 5 & low\;\, e\;\, s\;\, t \\
  \bottomrule
  \end{tabular}
  \end{center}

  \textbf{Round 3}: (e, s) appears 5 times---merge into ``es'':
  lowest $\rightarrow$ low\;\, es\;\, t.

  Continuing this process, we accumulate merge rules: (l, o), (lo, w),
  (e, s), \ldots\ The final vocabulary consists of all original characters
  plus every merged token.
\end{keyconcept}

\begin{pseudocode}[BPE Training]
def train_bpe(corpus: str, num_merges: int) -> tuple[list[tuple[str, str]], dict[str, int]]:
    word_freqs = pretokenize(corpus)                # {"lowest": 5, "low": 2, ...}
    # Split each word into characters
    splits: dict[str, list[str]] = {
        word: list(word) for word in word_freqs     # "lowest" -> ["l", "o", "w", "e", "s", "t"]
    }
    merges: list[tuple[str, str]] = []
    for i in range(num_merges):
        pair_counts = count_pairs(splits, word_freqs) # all adjacent pairs
        best_pair = max(pair_counts, key=pair_counts.get) # most frequent pair
        merges.append(best_pair)
        splits = merge_pair(splits, best_pair)      # replace (a, b) with "ab" everywhere
    vocab = build_vocab(splits, merges)             # chars + merged tokens + special tokens
    return merges, vocab
\end{pseudocode}

\subsection{BPE Encoding}

Once the merge rules are learned, encoding new text is deterministic:
split each word into characters, then replay the merge rules in order.

\begin{pseudocode}[BPE Encoding]
def encode_bpe(text: str, merges: list[tuple[str, str]], vocab: dict[str, int]) -> list[int]:
    words = pretokenize(text)                       # regex-based pre-tokenization
    token_ids: list[int] = []
    for word in words:
        symbols = list(word)                        # split into characters
        for (a, b) in merges:                        # replay merges in training order
            symbols = merge_pair_in_list(symbols, a, b)
        for symbol in symbols:
            # Note: for byte-level BPE this fallback never triggers,
            # since every byte is in the base vocabulary.
            token_ids.append(vocab.get(symbol, vocab["<|unk|>"]))
    return token_ids
\end{pseudocode}

\begin{keyconcept}
  \textbf{BPE key insight}: merge rules are ordered---encoding must
  replay them in the \emph{same order} as training. If we merged (l, o)
  before any other pair, then during encoding we must also apply that
  merge first. If we applied merges out of order, we could get different
  tokenizations and the model wouldn't recognize the resulting tokens.

  GPT-2's tokenizer uses byte-level BPE. It first splits text with a
  regex that separates contractions, numbers, and punctuation into
  reasonable chunks, then applies BPE merges learned from a large corpus.
  Its vocabulary contains 50,257 tokens.
\end{keyconcept}

\begin{keyconcept}[Encoding Efficiency]
  The \texttt{encode\_bpe} pseudocode above iterates over all merge rules
  for every word, giving a cost of $O(M \times L)$ per word where $M$ is
  the number of merges and $L$ is the word length. Real implementations
  (including GPT-2's) are more efficient: instead of scanning all merges,
  they repeatedly find the highest-priority merge \emph{currently present}
  in the word using a priority queue, reducing the cost to $O(L^2 \log L)$.
  The simple replay approach is pedagogically clear but not suitable for
  production use.
\end{keyconcept}

\begin{keyconcept}[Choosing Vocabulary Size]
  How many merges should you perform? This is controlled by the
  \texttt{num\_merges} hyperparameter. Typical values range from 30{,}000 to
  100{,}000. More merges means a larger vocabulary (so shorter sequences)
  but also more memory for the embedding table and a larger softmax at the
  output layer. Fewer merges means longer sequences (more compute for
  attention) but a smaller model. After training, the tokenizer is saved
  as two files: the merge rules list (ordered) and the vocabulary mapping
  (token $\rightarrow$ ID). To tokenize new text, you load both files and
  apply the \texttt{encode\_bpe} procedure.
\end{keyconcept}

\section{Next-Token Prediction Dataset}

Language models are trained via \emph{next-token prediction}: given a
sequence of tokens, predict the next one. To create training examples
from a long tokenized corpus, we use a sliding window.

\begin{pseudocode}[Sliding-Window Dataset]
class SlidingWindowDataset:
    def __init__(self, token_ids: list[int], max_length: int, stride: int):
        self.token_ids = token_ids                  # full corpus as IDs
        self.max_length = max_length                # context window size
        self.stride = stride                        # step between windows

    def __getitem__(self, idx: int) -> tuple[Tensor, Tensor]:
        start = idx * self.stride
        end = start + self.max_length
        input_ids = torch.tensor(self.token_ids[start:end])       # [max_length]
        target_ids = torch.tensor(self.token_ids[start+1:end+1])  # [max_length]
        return input_ids, target_ids

    def __len__(self) -> int:
        return (len(self.token_ids) - self.max_length) // self.stride
        # Ensures end+1 <= len(token_ids), so no out-of-bounds access
\end{pseudocode}

\begin{keyconcept}
  The target is the input shifted by one position: if the input is
  ``The cat sat'', the target is ``cat sat on''. This way, at each
  position the model predicts the \emph{next} token given everything
  before it.

  The \textbf{stride} controls overlap between windows. Smaller stride
  $=$ more training examples but more redundancy (adjacent windows share
  most of their tokens). When $\text{stride} = \text{max\_length}$,
  there is no overlap and each token appears in exactly one example.
\end{keyconcept}

\section*{Review Questions}

\begin{reviewquestion}[1]
  Why does character-level tokenization produce much longer sequences than
  word-level tokenization? Estimate the approximate length ratio for a
  typical English sentence.
\end{reviewquestion}

\begin{reviewquestion}[2]
  Why does word-level tokenization suffer from vocabulary explosion, and
  how does BPE address this?
\end{reviewquestion}

\begin{reviewquestion}[3]
  Give an example of a word that would be tokenized differently under
  word-level vs.\ subword tokenization. Why does the subword approach
  avoid the OOV problem?
\end{reviewquestion}

\begin{reviewquestion}[4]
  In the SimpleTokenizer, what happens when a word is not in the
  vocabulary? How does this compare to what happens with a BPE tokenizer
  encountering a rare word?
\end{reviewquestion}

\begin{reviewquestion}[5]
  What is the role of the \texttt{<|eos|>} token? Why would you insert it
  between documents in a training corpus?
\end{reviewquestion}

\begin{reviewquestion}[6]
  In BPE training, suppose the pair (e, r) is the most frequent across the
  corpus. What happens after this pair is merged? Does it affect the
  frequency of other pairs?
\end{reviewquestion}

\begin{reviewquestion}[7]
  In the worked example above, what would the tokens for ``lowest'' look
  like after 4 rounds of merging?
\end{reviewquestion}

\begin{reviewquestion}[8]
  Why must BPE merge rules be replayed in the same order during encoding
  as they were learned during training? Construct a small example where
  applying merges in a different order produces a different tokenization.
\end{reviewquestion}

\begin{reviewquestion}[9]
  How does the stride parameter in the sliding-window dataset affect both
  the number of training examples and their redundancy? What are the
  tradeoffs of using a very small stride vs.\ setting stride equal to
  max\_length?
\end{reviewquestion}

\begin{reviewquestion}[10]
  GPT-2 uses byte-level BPE. What does ``byte-level'' mean, and why does
  it guarantee that any input string (including emojis, multilingual text,
  etc.) can be tokenized without an \texttt{<|unk|>} token?
\end{reviewquestion}

\begin{reviewquestion}[11]
  Suppose you have a corpus with the words ``running'' (freq 10) and
  ``runner'' (freq 8). After the first BPE merge creates the token ``run'',
  what does the tokenization of ``running'' and ``runner'' look like?
\end{reviewquestion}

\begin{reviewquestion}[12]
  Why is pre-tokenization (splitting text into words with a regex before
  BPE) useful? What would happen if you applied BPE directly on raw
  character sequences without pre-tokenization?
\end{reviewquestion}

\section*{Problems}

\begin{problem}[1]
  Implement a \texttt{SimpleTokenizer} class with \texttt{encode} and
  \texttt{decode} methods. Your tokenizer should split on whitespace,
  maintain a vocabulary dictionary, and handle unknown tokens with a
  special \texttt{<|unk|>} token.
\end{problem}

\begin{problem}[2]
  Implement BPE training from scratch. Given a small corpus (e.g., a few
  paragraphs of text), write a function that:
  \begin{enumerate}
    \item Pre-tokenizes the corpus into words with frequency counts
    \item Splits each word into characters
    \item Iteratively finds and merges the most frequent pair
    \item Returns the list of merge rules
  \end{enumerate}
  Test your implementation on the example from this chapter and verify
  that you get the same merge sequence.
\end{problem}

\begin{problem}[3]
  Implement the \texttt{encode} method for BPE using your trained merge
  rules. For each word: split into characters, then replay the merge
  rules in order. Ensure that your encoding reproduces the same
  tokenization as the training process. Test it by encoding ``lowest''
  and ``low'' and checking that the results match what you observed
  during training.
\end{problem}

\begin{problem}[4]
  Implement the \texttt{SlidingWindowDataset} class. Use a real
  corpus (e.g., a public domain book from Project Gutenberg), tokenize
  it with your BPE tokenizer, and create the dataset with
  \texttt{max\_length=256} and \texttt{stride=128}. Print out the first
  3 examples (both input and target) and verify that the target is the
  input shifted by one position.
\end{problem}

\begin{problem}[5]
  Investigate the effect of stride on dataset size and redundancy. Using
  your \texttt{SlidingWindowDataset} with \texttt{max\_length=256},
  compute the number of examples for stride values of 32, 64, 128, and
  256. For each setting, also estimate the percentage of token overlap
  between consecutive examples. What are the memory and training
  implications of each setting?
\end{problem}